\section{Proof of the symplectic invariance of \texorpdfstring{$F^{(0)}$}{F0} and \texorpdfstring{$F^{(1)}$}{F1}}
\label{appsymplinvF0F1}


\subsection{\texorpdfstring{$F^{(1)}$}{F1}}
Let us study how $F^{(1)}$ is changed under the exchange of the roles of $x$ and $y$.
For this purpose, we define the images of $F^{(1)}$ and $W_1^{(1)}(p)$ under this
transormation, i.e.
\beq
\widehat{W}_1^{(1)}(p):= - \Res_{q \to {\bf b}} {\int_{\tilde{q}}^q B(p,\xi) \over 2 (x(q)-x(\tilde{q})) dy(q)} B(q, \tilde{q})
\eeq
and
\beq
\widehat{F}^{(1)} = -{1 \over 2 } ln \tau_{By} - {1 \over 24} ln \prod_j x'(b_j)
\eeq
where ${\bf b}$ denotes the set of $y$-branch points,
$\tilde{q}$ is the only point satisfying $y(q) = y(\tilde{q})$ and approaching a branch point $b_j$
when $q \to b_j$, $\tau_{By}$ is the Bergmann tau function associated to $y$ and $x'(b_j)= {dx(b_j) \over d\tilde{z}_j(b_j)}$.
According to theorem \ref{variat}, for any variation $\Omega$ of the curve, the variation of the free energy reads
\beq
\delta_\Omega \widehat{F}^{(1)} = \int_{\cal{C}} \widehat{W}_1^{(1)}(p) \L(p).
\eeq
Thus, the variation of the difference between the two "free energies" reads
\beq
\delta_\Omega \left( F^{(1)} - \widehat{F}^{(1)} \right) = \int_{\cal{C}} \left( W_1^{(1)}(p) + \widehat{W}_1^{(1)}(p) \right) \L(p) .
\eeq

In order to evaluate this quantity, one needs the following lemma:
\bl
For any choice of variable $z$:
\beq\label{sumW}
W_1^{(1)}(p) + \widehat{W}_1^{(1)}(p) = {1 \over 24} d_p \left[{1 \over x' y'} \left( 2  S_{Bz}(p)
+ {x'' y'' \over x' y'} + {x''^2 \over x'^2} - {x''' \over x'} + {y''^2 \over y'^2} - {y''' \over y'} \right) \right]
\eeq
where the derivatives are taken with respect to the variable $z$ and $S_{Bz}$ denotes
the Bergmann projective connection associated to $z$.
\el

\proof{
From the definitions , one easily derives that
\beq\label{sumW1}
W_1^{(1)}(p) + \widehat{W}_1^{(1)}(p) = \Res_{q \to p} {B(p,q) \over (x(p)-x(q)) (y(p)-y(q))}.
\eeq

Let us now expand the integrand in terms of an arbitrary local variable $z$ when $p \to q$. The different factors
read
\beq
{B(p,q) \over dz(p) dz(q)}= {1 \over (z(p)-z(q))^2} + {1 \over 6} S_{Bz}(p) - {z(p)-z(q) \over 12} S_{Bz}'(p)+ O((z(p)-z(q))^2)
\eeq
and
\bea
{1 \over y(p)-y(q)} &=& {1 \over (z(p)-z(q)) y'(p)} \Big[ 1 - (z(p)-z(q)) {y'' \over 2 y'} \cr 
&& \; \; + (z(p)-z(q))^2 \left( {y''^2 \over 4 y'^2} - {y'''\over 6 y'}\right) \cr
&& \; \; \; \;   + {(z(p)-z(q))^3 \over 24} \left(- {y'''' \over y'} +4 {y''' y'' \over y'^2} - {3 y''^3 \over y'^3} \right) \cr
&& \qquad + O((z(p)-z(q))^4) \Big] .\cr
\eea
Inserting it altogether inside \eq{sumW1}, one can explicitly compute the residues and one recognizes the
formula \ref{sumW}.
}

We can now prove the following theorem stating the symmetry of $F^{(1)}$ under the exchange of $x \leftrightarrow y$:
\bl\label{symF1}
The two free energies transform in the same way under any variation of the moduli of the curve:
\beq
\delta_\Omega  F^{(1)} = \delta_\Omega \widehat{F}^{(1)}
\eeq
\el

\proof{
We already know that
\beq
\delta_\Omega \left( F^{(1)} - \widehat{F}^{(1)} \right) = \int_{p \in \cal{C}}
{1 \over 24} d_p \left[{1 \over x' y'} \left( 2  S_{Bz}(p)
+ {x'' y'' \over x' y'} + {x''^2 \over x'^2} - {x''' \over x'} + {y''^2 \over y'^2} - {y''' \over y'} \right) \right]
\L(p)
\eeq
for an arbitrary variable $z$.

We just have to check that this quantity vanishes for the transformations corresponding to varying the moduli
of the curve. Because the function inside the differential wrt $p$ vanishes at the poles of $ydx$, one can
check that it is indeed the case.
}


Thus the first correction to the free energy satisfies the following variation under symplectic transformations:
\bt
$F^{(1)}$ does not change under the following transformations of $\curve$:

$\bullet$ \vspace{0.3mm} $y\to y +R(x)$ where $R$ is a rational function.

$\bullet$ \vspace{0.3mm} $y\to c y$ and $x\to {1\over c}x$ where $c$ is a non-zero complex number.

$\bullet$ \vspace{0.3mm} $x\to x +R(y)$ where $R$ is a rational function.

$\bullet$ \vspace{0.3mm} $y\to x$ and $x\to y$.

$F^{(1)}$ is shifted by a multiple of $ i \pi/12$ when $\curve$ is changed by $x \to - x$.

\et

\proof{The first transformation is obvious from the definition since neither $ln \tau_{Bx}$ nor $y'(a_i)$ changes.

The second one follows from theorem 2 in \cite{EKK} for $f=x$ and $g={x \over c}$ which shows that
the variations of $ln(\tau_{Bx})$ and $y'(a_i)$ compensate.

The fourth one is nothing but lemma (\ref{symF1}) and it gives the third one when combined with the first.

The last transformation holds because $ln \tau_{Bx}$ is left unchanged and $y'(a_i)$ changes sign.

}

\subsection{\texorpdfstring{$F^{(0)}$}{F0}}

\bt
$F^{(0)}$ does not change under the following transformations of $\curve$:

$\bullet$ \vspace{0.3mm} $y\to y +{\cal{P}}(x)$ where ${\cal{P}}$ is a polynomial.

$\bullet$ \vspace{0.3mm} $y \to y$ and $x \to -x$.

$\bullet$ \vspace{0.3mm} $y\to c y$ and $x\to {1\over c}x$ where $c$ is a non-zero complex number.

$\bullet$ \vspace{0.3mm} $x\to x +{\cal{P}}(y)$ where ${\cal{P}}$ is a polynomial.

$\bullet$ \vspace{0.3mm} $y\to x$ and $x\to y$.

\et

\proof{
The second and third transformations come from the structure of $F^{(0)}$ which is bilinear in objects proportionnal to $ydx$.

One can obtain the last one following the same method as for $F^{(1)}$: we compute the variation of the difference
of the two free energies under the changes of the moduli of $\curve$ and show that they vanish.

Let us now show the first invariance. In this case, the variation of the free energy is
$\delta_{{\cal{P}}(x)dx} F^{(0)}$. Since ${\cal{P}}(x)$ is a polynomial one can write it
\beq
{\cal{P}}(x(p)) = \Res_{q \to p} B(p,q) {\cal{Q}}(x(q)) = - \sum_\alpha \Res_{q \to \alpha} B(p,q) {\cal{Q}}(x(q)),
\eeq
where ${\cal{Q}}(x)$ is also a polynomial in $x$. Then, theorem (\ref{variat}) implies the invariance
of $F^{(0)}$ under this transformation.

The fourth transformation is a combination of the first and the last one.


}




